\chapter{Practical Implementation of the Project}
\label{ch:implementation}

% Figures are non-floating (center + \captionof) so they stay in place.
% Requires in the main preamble: \usepackage{caption}  % or capt-of

This chapter describes how we implemented \texttt{roughapp}. We first summarise
the system requirements and then explain the architecture and the data model.
Finally, we outline the technology stack and present the user interface.

The application is a web tool for analysing inconsistency in decision tables and
for comparing inconsistency-resolution strategies. A user uploads a CSV dataset,
selects the decision attribute, inspects rough-set measures (in particular
$\gamma$ and $\alpha$) and conflicting equivalence classes, and then applies one
or more strategies (MCD, qualitative removal, MVD, GD). Outputs are stored as
new datasets so they can be analysed again, exported, or used as input for
further experiments. Decision trees (CART) are built on demand to compare model
complexity before and after resolution.

\section{Requirements}

\subsection{Functional requirements}

\begin{description}

\item[FR-1 (Accounts and anonymous work).]
A user can create an account with an e-mail address and a password and log in
with the same credentials. The system can also be used without an account: a
visitor works in a temporary session. If the visitor registers or logs in while
such a session is still active, the session data is transferred into the account.

\item[FR-2 (Ownership-based access).]
Built-in datasets are public. Datasets, experiments, and decision trees created
by a registered user or an anonymous session are visible only to their owner.

\item[FR-3 (Dataset management).]
A user uploads a decision table as a CSV file and selects the decision attribute.
The system stores the dataset name, dimensions, and the chosen decision attribute.
Uploaded and built-in datasets can be listed and previewed page by page.

The CSV may include an optional \texttt{type row} immediately below the header,
containing one marker per column:
\texttt{c} (categorical), \texttt{l} (logical/binary), \texttt{n} (numeric),
\texttt{i} (integer), and \texttt{r} (real). If the type row is absent, all
attributes are treated as categorical. The type row is not stored separately in
the database: it remains part of the CSV file and is written automatically when
the system saves derived or resolved tables.

\item[FR-4 (Inconsistency analysis).]
For a chosen dataset, the system computes equivalence classes, lower and upper
approximations, the boundary region, $\alpha$ and $\gamma$, and lists conflicting
classes with example objects. It can also rank condition attributes by importance,
measured as the drop in $\gamma$ after removing each attribute.

\item[FR-5 (Introducing inconsistency).]
An inconsistent table can be derived from a consistent one by duplicating objects
with a flipped decision, removing selected attributes, or repeating duplication
until $\gamma$ reaches a target. The result becomes a new dataset linked to its
source.

\item[FR-6 (Handling missing values).]
Missing values are detected and counted per column. The user can drop incomplete
rows or fill the gaps using the most common value for nominal attributes and the
mean for numeric attributes, either globally or within each decision class. The
completed table becomes a new dataset linked to its source.

\item[FR-7 (Running experiments).]
One or more strategies are applied to a dataset in a single experiment, with
tie-breaking policies and a random seed configurable. For every strategy, the
system records $\gamma$, $\alpha$, row counts, running time, and details of how
conflicts were handled. The input dataset is never modified; each output is
stored as a new dataset.

\item[FR-8 (Decision tree comparison).]
A decision tree (currently CART) can be built on demand for any dataset (original
or resolved). The system reports depth, average leaf depth, and node count.
Numeric markers (\texttt{n}, \texttt{i}, \texttt{r}) are encoded as numeric
features, while categorical and logical markers (\texttt{c}, \texttt{l}) are
one-hot encoded. A comparison view places previously built trees side by side.

\item[FR-9 (Export).]
Any resolved table can be downloaded as a CSV file and used in external tools.

\end{description}

Figure~\ref{fig:usecase} summarises these requirements. A single actor, the
user, covers both anonymous work in a temporary session and work after
registration or login (FR-1, FR-2). Guest data can be adopted into an account when
the user registers or logs in.

\begin{center}
\includegraphics[width=0.9\textwidth]{figures/use_case.png}
\captionof{figure}[Use case diagram]{Use case diagram of the system.\\{\footnotesize Source: Own work}}
\label{fig:usecase}
\end{center}

Figure~\ref{fig:activity_workflow} shows a typical analytical path: choose or
upload a dataset, optionally analyse it, complete missing values or introduce
inconsistency when needed, run an experiment, inspect results, and export or
reuse resolved tables.

\begin{center}
\includegraphics[width=0.85\textwidth]{figures/activity_workflow.png}
\captionof{figure}[Main analytical workflow]{Main analytical workflow.\\{\footnotesize Source: Own work}}
\label{fig:activity_workflow}
\end{center}

\subsection{Non-functional requirements}

\begin{description}

\item[NFR-1 (Reproducibility).]
Every operation with a random component takes an explicit seed, and all
experiment parameters are stored. Running the same operation on the same data
with the same seed produces the same result.

\item[NFR-2 (Correctness and testability).]
The rough-set measures and all four resolution strategies are covered by unit
tests and must reproduce the worked examples from Chapter~4 exactly.

\item[NFR-3 (Portability and deployment).]
The system runs from one codebase in two configurations: locally with an embedded
database and as a containerised stack with a client--server database. Switching
between them requires no code changes, and the complete stack starts with a
single command.

\item[NFR-4 (Extensibility).]
Adding a resolution or introduction strategy requires one new module and one
registry entry; existing strategies and API code do not need to change.

\item[NFR-5 (Security).]
Passwords are never stored or logged in plain text. Every operation on a private
resource checks ownership. In production, only the web frontend is externally
accessible.

\item[NFR-6 (Interoperability).]
Input and output use plain CSV files.

\item[NFR-7 (Usability).]
The main analytical functions work without an account, and results are displayed
without full-page reloads.

\item[NFR-8 (Responsiveness under load).]
Long-running experiments, tree generation and attribute-importance calculations
must not prevent the server from responding to other requests.

\end{description}

\section{System Architecture}

The system follows a three-tier client--server architecture: a browser frontend,
an application server, and a relational database. The frontend is a single-page
application that communicates with the server through a JSON REST API. In
production, only the frontend is publicly accessible; the API and database remain
on an internal network.

Internally, the application server is organised into three layers:
\begin{itemize}
\item \textbf{Presentation layer} (\texttt{api/}): HTTP endpoints, request parsing,
validation, and response formatting.
\item \textbf{Business layer} (\texttt{business/}): services, access rules,
algorithms, and transaction logic. Algorithms are exposed through a registry and
common interfaces, which makes adding new strategies straightforward (NFR-4).
\item \textbf{Data access layer} (\texttt{db/}): database queries and file storage
logic.
\end{itemize}

Figure~\ref{fig:arch_layers} shows how these layers sit between the browser and
the stored data.

\begin{center}
\includegraphics[width=0.85\textwidth]{figures/component_architecture.png}
\captionof{figure}[Layered architecture]{Layered architecture of the application.\\{\footnotesize Source: Own work}}
\label{fig:arch_layers}
\end{center}

The independent \texttt{core/} module contains rough-set operations (equivalence
classes, approximations, $\alpha$, $\gamma$) and missing-value detection. It does
not depend on the web framework or the database, which makes it easy to test in
isolation (NFR-2).

Structured metadata is stored in PostgreSQL in production and in SQLite locally
(NFR-3). Decision tables are stored as CSV files, while the database stores their
metadata and file paths. Files are written before the database commit and removed
if the transaction fails; the database record remains the source of truth.

Heavy computations run in worker threads so that the server can continue handling
requests while experiments or tree generation are running (NFR-8).

\section{Data Model}

The database schema is managed through Alembic migrations, which run automatically
when the server starts. The same migrations are used for SQLite and PostgreSQL, so
both configurations share the same structure.

For readability, the schema is presented in five areas: users and ownership,
dataset management, experiment configuration, experiment execution and results,
and decision-tree generation. Reference tables are shown separately.

\subsection{Users and ownership}

Data may belong to a registered \texttt{user} or to an anonymous
\texttt{guest\_session}, identified by a hashed cookie token with a limited
lifetime. Datasets, experiments, and tree runs store either a \texttt{user\_id}
or a \texttt{guest\_session\_id}. Built-in public datasets may have neither,
while every experiment must have exactly one owner.

\begin{center}
\includegraphics[width=0.9\textwidth]{figures/users_logical.png}
\captionof{figure}[Users and ownership --- logical model]{Users and ownership --- logical model.\\{\footnotesize Source: Own work}}
\label{fig:log_users}
\end{center}

\begin{center}
\includegraphics[width=0.9\textwidth]{figures/users_relational.png}
\captionof{figure}[Users and ownership --- relational model]{Users and ownership --- relational model.\\{\footnotesize Source: Own work}}
\label{fig:erd_users}
\end{center}

\subsection{Dataset management}

Each \texttt{dataset} stores basic metadata about a decision table: its name, CSV
path, dimensions, the decision attribute, and cached values of $\gamma$ and
$\alpha$. The dataset origin is stored in \texttt{dataset\_provenance} as an
upload, a derived dataset, or an experiment output, with an optional reference to
the source dataset. This makes every generated table a regular dataset that can be
analysed again and reused.

If the source dataset is deleted, derived datasets remain available; only their
\texttt{source\_dataset\_id} is set to \textsc{null}.

\begin{center}
\includegraphics[width=0.9\textwidth]{figures/datasets_logical.png}
\captionof{figure}[Dataset management --- logical model]{Dataset management --- logical model.\\{\footnotesize Source: Own work}}
\label{fig:log_datasets}
\end{center}

\begin{center}
\includegraphics[width=0.9\textwidth]{figures/datasets_relational.png}
\captionof{figure}[Dataset management --- relational model]{Dataset management --- relational model.\\{\footnotesize Source: Own work}}
\label{fig:erd_datasets}
\end{center}

\subsection{Experiment configuration}

An \texttt{experiment} stores the settings of a run: the input dataset, random
seed, tie-breaking policies, and owner. It does not store measurements. Input
metrics are taken from the dataset, while outcomes are stored with experiment
results.

\begin{center}
\includegraphics[width=0.9\textwidth]{figures/experiment_config_logical.png}
\captionof{figure}[Experiment configuration --- logical model]{Experiment configuration --- logical model.\\{\footnotesize Source: Own work}}
\label{fig:log_expconfig}
\end{center}

\begin{center}
\includegraphics[width=0.9\textwidth]{figures/experiment_config_relational.png}
\captionof{figure}[Experiment configuration --- relational model]{Experiment configuration --- relational model.\\{\footnotesize Source: Own work}}
\label{fig:erd_expconfig}
\end{center}

\subsection{Experiment execution and results}

Each strategy run creates one \texttt{experiment\_result} storing the strategy,
its runtime, and a JSON \texttt{details} field describing how conflicts were
handled. Numeric values (such as \texttt{gamma}, \texttt{alpha}, row counts, and
strategy-specific measures) are stored in \texttt{experiment\_metrics} as
name--value pairs. This design allows adding new metrics without schema changes.

Each result may point to one output dataset through
\texttt{output\_dataset\_id}. The link is optional: if the resolved CSV is deleted
later, the metrics remain and the foreign key is cleared. Because outputs are
ordinary datasets, a resolved table can go straight into another experiment or
into tree building.

\begin{center}
\includegraphics[width=0.95\textwidth]{figures/experiment_results_logical.png}
\captionof{figure}[Experiment execution and results --- logical model]{Experiment execution and results --- logical model.\\{\footnotesize Source: Own work}}
\label{fig:log_expresults}
\end{center}

\begin{center}
\includegraphics[width=0.95\textwidth]{figures/experiment_results_relational.png}
\captionof{figure}[Experiment execution and results --- relational model]{Experiment execution and results --- relational model.\\{\footnotesize Source: Own work}}
\label{fig:erd_expresults}
\end{center}

\subsection{Decision-tree generation}

A \texttt{tree\_run} records that a decision tree was built for a given dataset
with a chosen algorithm (currently CART). Structural measures such as depth and
node count are stored in \texttt{tree\_metrics}. Because resolved experiment
outputs are saved as datasets, comparing trees before and after resolution is
simply comparing runs on different datasets.

\begin{center}
\includegraphics[width=0.9\textwidth]{figures/trees_logical.png}
\captionof{figure}[Decision-tree generation --- logical model]{Decision-tree generation --- logical model.\\{\footnotesize Source: Own work}}
\label{fig:log_trees}
\end{center}

\begin{center}
\includegraphics[width=0.9\textwidth]{figures/trees_relational.png}
\captionof{figure}[Decision-tree generation --- relational model]{Decision-tree generation --- relational model.\\{\footnotesize Source: Own work}}
\label{fig:erd_trees}
\end{center}

\subsection{Reference tables}

Strategy, policy, and algorithm names are stored in lookup tables rather than as
free text. This makes validation simpler and keeps the database consistent.
Figure~\ref{fig:erd_reference} shows the reference tables: four resolution
strategies, the two tie-policy lists, and the tree-algorithm list.

\begin{center}
\includegraphics[width=0.8\textwidth]{figures/reference_tables.png}
\captionof{figure}[Reference tables]{Reference tables.\\{\footnotesize Source: Own work}}
\label{fig:erd_reference}
\end{center}

\subsection{Integrity rules and deletion}

Foreign keys and unique constraints keep the schema consistent. Each experiment
can have at most one result per strategy, each result can have at most one metric
with a given name, each result can point to at most one output dataset, and each
owner can have at most one dataset with the same \texttt{storage\_path}. A small
number of \texttt{CHECK} constraints enforce valid ownership types (registered
user, guest, or built-in). Other rules, such as valid strategy and policy names,
are enforced in the API.

Deleting a dataset removes its CSV file and any trees built from it after the
database transaction is committed. Derived datasets and experiment outputs remain,
but their \texttt{source\_dataset\_id} is set to \textsc{null}. If the dataset was
used as input for experiments, the API first displays affected experiments and
deletes their records only after confirmation; resolved CSV outputs remain
available as separate datasets.

\subsection{Normal Forms}

The schema satisfies the first three normal forms. Each table has a primary key,
columns contain atomic values, and repeating data (metrics) is stored in child
tables (1NF). Non-key attributes depend on the whole key (2NF). Lookup tables
prevent transitive dependencies (3NF).

Two deliberate exceptions remain. The JSON \texttt{details} field stores metadata
that is naturally processed as a whole. The \texttt{dataset} table also caches
measured metric values to avoid reading every CSV file when displaying dataset
lists. Since dataset files are immutable, cached values remain consistent.

\section{Used Technologies}
\label{sec:tech}

\section{User Interface}
\label{sec:ui}
